\documentclass[aspectratio=169,12pt]{beamer}
\usepackage{fontspec}
\usepackage[portuguese]{babel}
\usepackage{xcolor,listings,tikz,booktabs,amsmath,amssymb,graphicx,multicol,array,colortbl}
\usetikzlibrary{arrows.meta,positioning,shapes.geometric,calc}

\definecolor{navy}{HTML}{0D1B3E}
\definecolor{gold}{HTML}{F5A623}
\definecolor{qgreen}{HTML}{1E8B4C}
\definecolor{qblue}{HTML}{1A6EAE}
\definecolor{codebg}{HTML}{EEF3F8}

\usetheme{default}
\useinnertheme{circles}
\setbeamertemplate{navigation symbols}{}
\setbeamertemplate{footline}{%
  \hfill{\color{gray!60}\scriptsize\insertframenumber/\inserttotalframenumber}\hspace{.6em}\vspace{.4em}}

\setbeamercolor{frametitle}{bg=navy,fg=white}
\setbeamercolor{structure}{fg=navy}
\setbeamercolor{alerted text}{fg=gold}
\setbeamercolor{item}{fg=gold}
\setbeamercolor{subitem}{fg=navy}
\setbeamercolor{block title}{bg=qblue,fg=white}
\setbeamercolor{block body}{bg=qblue!8,fg=black}
\setbeamercolor{block title alerted}{bg=gold,fg=navy}
\setbeamercolor{block body alerted}{bg=gold!15,fg=navy}
\setbeamercolor{block title example}{bg=qgreen,fg=white}
\setbeamercolor{block body example}{bg=qgreen!10,fg=black}
\setbeamerfont{frametitle}{size=\large,series=\bfseries}
\setbeamerfont{block title}{series=\bfseries}

\setbeamertemplate{title page}{%
  \begin{tikzpicture}[remember picture,overlay]
    \fill[navy](current page.north west)rectangle(current page.south east);
    \fill[gold]([yshift=.7cm]current page.south west)rectangle(current page.south east);
    \node[anchor=south,yshift=.73cm,navy,font=\tiny\bfseries]at(current page.south)
      {INTENSIVÃO QUANT AI \textbullet{} IMPACT UFSCAR};
  \end{tikzpicture}
  \vspace{.65cm}
  \begin{center}
    {\color{gold}\scriptsize\bfseries INTENSIVÃO QUANT AI \textbullet{} IMPACT UFSCAR}\\[.4cm]
    {\color{white}\LARGE\bfseries\inserttitle}\\[.25cm]
    {\color{gold!85}\normalsize\insertsubtitle}\\[.5cm]
    {\color{white!70}\small\insertauthor}\\[.1cm]
    {\color{white!50}\footnotesize\insertdate}
  \end{center}}

\lstdefinestyle{python}{language=Python,
  basicstyle=\ttfamily\scriptsize\color{qblue},
  keywordstyle=\color{navy}\bfseries,
  stringstyle=\color{qgreen!80!black},
  commentstyle=\color{gray!70}\itshape,
  backgroundcolor=\color{codebg},
  frame=l,framerule=2pt,rulecolor=\color{gold},
  xleftmargin=1em,breaklines=true,showstringspaces=false,tabsize=4,
  extendedchars=false,inputencoding=ascii}
\lstset{style=python}

\author{Felipe Maldonado\\{\scriptsize VP Diretoria Quant~\textbullet{}~Impact UFSCAR}}
\date{Intensivão Quant AI 2026}
\title{Aula 04 --- Sinal v1}
\subtitle{Construindo o Sinal de Momentum Cross-Sectional}

\begin{document}

\begin{frame}[plain]
\titlepage
\end{frame}

\begin{frame}[fragile]{Nesta Aula}
\begin{columns}[T]
  \column{0.5\textwidth}
\begin{block}{Teoria (20 min)}
\begin{itemize}
  \item O pipeline quant de geração de alfa
  \item Anatomia do sinal 12-1 de momentum
  \item Information Coefficient (IC): o que é e por que importa
  \item IC ratio e IR: avaliando poder preditivo
\end{itemize}
\end{block}
  \column{0.47\textwidth}
\begin{block}{Código (40 min)}
\begin{itemize}
  \item \texttt{calcular\_sinal\_v1()} — implementação do sinal 12-1
  \item Ranking cross-sectional mensal
  \item \texttt{calcular\_ic\_serie()} — Spearman mês a mês
  \item Visualização: IC barras e rolling IC
  \item Salvar \texttt{sinal\_v1.parquet}
\end{itemize}
\end{block}
\end{columns}
\end{frame}

\begin{frame}[fragile]{O Fator Momentum — Sinal 12-1}
\begin{columns}[T]
  \column{0.5\textwidth}
\textbf{Definição formal:}
\[
\text{Sinal}_{i,t} = \sum_{k=2}^{12} r_{i,t-k}
\]
\vspace{.1em}
Em pandas:
\begin{lstlisting}
ret.shift(2).rolling(11).sum()
\end{lstlisting}
\vspace{.3em}
\begin{itemize}
  \item \texttt{shift(2)}: exclui o mês mais recente (reversão de curto prazo: Jegadeesh, 1990)
  \item \texttt{rolling(11)}: soma dos 11 meses anteriores ao mês exluído
  \item Total: janela de formação de 12 meses, skip 1 mês
\end{itemize}
  \column{0.47\textwidth}
\begin{block}{Intuição Econômica}
\begin{itemize}
  \item Underreaction de investidores a notícias (Hong \& Stein, 1999)
  \item Herding: fundos seguem vencedores recentes
  \item Momentum institucional: trimestral window dressing
  \item Prêmio por risco de crash (Daniel \& Moskowitz, 2016)
\end{itemize}
\end{block}\vspace{.3em}
\begin{block}{Ranking Cross-Sectional}

Cada mês $t$: ranquear todos os ativos pelo sinal.\\
Comprar \textbf{top 10\%} (quintil superior).\\
Portfólio long-only, equal-weight entre selecionados.
\end{block}
\end{columns}
\end{frame}

\begin{frame}[fragile]{Information Coefficient (IC)}
\begin{columns}[T]
  \column{0.5\textwidth}
\textbf{Definição:}
\[ \text{IC}_t = \rho_S\!\left(\text{Sinal}_{i,t},\; r_{i,t+1}\right) \]
onde $\rho_S$ é a correlação de Spearman.
\vspace{.4em}
\begin{itemize}
  \item $\text{IC} > 0$: sinal tem poder preditivo
  \item $\text{IC} = 0$: sinal aleatório
  \item Referência: IC médio $> 0{,}05$ é considerado bom
  \item Spearman (vs Pearson): robusto a outliers de retorno
\end{itemize}
\vspace{.3em}
\begin{alertblock}{IC vs t-statistic}
IC médio isolado é insuficiente. Calcule $t = \frac{\overline{IC}}{\sigma_{IC}/\sqrt{T}}$.
Um IC de 0{,}04 com $T=60$ meses dá $t \approx 2{,}0$ — marginalmente significativo.
\end{alertblock}
  \column{0.47\textwidth}
\begin{block}{IR — Information Ratio}

\[ \text{IR} = \frac{\overline{IC}}{\sigma_{IC}} \]
\begin{itemize}
  \item Análogo ao Sharpe do sinal
  \item IR $> 0{,}5$: sinal de qualidade institucional
  \item Fundamento da Fundamental Law of Active Management (Grinold, 1989):
\end{itemize}
\[ \text{IR} \approx \text{IC} \times \sqrt{N} \]
onde $N$ = número de apostas independentes por período.
\end{block}
\end{columns}
\end{frame}

\begin{frame}[fragile]{O que Vamos Construir}
\begin{columns}[T]
  \column{0.52\textwidth}
\begin{block}{Funções a implementar}
\begin{itemize}
  \item \texttt{calcular\_sinal\_v1(ret\_mensais)} $\to$ DataFrame de sinais
  \item \texttt{rankear\_sinal(sinal)} $\to$ ranks normalizados 0--1
  \item \texttt{calcular\_ic\_serie(sinal, ret)} $\to$ Series de IC mensais
  \item \texttt{plot\_ic\_barras(ic\_serie)} $\to$ gráfico de barras com IC rolling
\end{itemize}
\end{block}
  \column{0.45\textwidth}
\begin{block}{Entradas (parquets)}
\begin{itemize}
  \item \texttt{retornos\_mensais\_limpo.parquet}
\end{itemize}
\end{block}
\vspace{.4em}
\begin{exampleblock}{Saídas (parquets)}
\begin{itemize}
  \item \texttt{sinal\_v1.parquet}
  \item Gráfico IC mensal e rolling
\end{itemize}
\end{exampleblock}
\end{columns}
\end{frame}

\begin{frame}[fragile]{Referências}
\begin{multicols}{2}
\tiny
\begin{itemize}
  \item \textbf{Jegadeesh, N. \& Titman, S.} (1993). Returns to buying winners and selling losers. \textit{Journal of Finance}, 48(1), 65--91.
  \item \textbf{Jegadeesh, N.} (1990). Evidence of predictable behavior of security returns. \textit{Journal of Finance}, 45(3), 881--898.
  \item \textbf{Grinold, R.} (1989). The fundamental law of active management. \textit{Journal of Portfolio Management}, 15(3), 30--37.
  \item \textbf{Hong, H. \& Stein, J.} (1999). A unified theory of underreaction, momentum trading and overreaction. \textit{Journal of Finance}, 54(6), 2143--2184.
  \item \textbf{Daniel, K. \& Moskowitz, T.} (2016). Momentum crashes. \textit{JFE}, 122(2), 221--247.
\end{itemize}
\end{multicols}
\end{frame}


\end{document}
